\scanpage{31}
\deleted{d'où une opération}
\[
\deleted{\psi:\check{\mathbb V}(L)\longrightarrow\underline{\operatorname{Autom}}\ \underline{\text{d'ex}}\,
(0\mathbin{-}L^{\otimes2}\mathbin{-}L\otimes\check B\mathbin{-}L\mathbin{-}0)}
\]
\deleted{et on a un}

scindage de \deleted{\unsure{l'}} l'extension
\[
0\mathbin{-}L\mathbin{-}\check B\mathbin{-}\underline O_S\mathbin{-}0 \tag*{(37)}
\]
qui décrit le fibré affine $E$, i.e. comme
un iso
\[
\check B\simeq\deleted{L\oplus}\underline O_S\oplus\underline{L}
\]
on aura un isom d'\uline{extension}
\[
\widetilde\alpha_s:\check B\otimes L\simeq L\oplus L^{\otimes2} \tag*{(38)}
\]
On trouve que \ldots

\noindent\textemdash\ \textemdash\ \textemdash

et on \deleted{\unsure{a}} trouve un iso. canonique
entre \unsure{l'groupe|le groupe} des translations de l'espace
affine $\Sigma^2(F)$ \deleted{\unsure{(de}}
\[
\operatorname{Trans}\Sigma^2(F)=F\mathop{\wedge}_{\check{\mathbb V}(L),\varphi_2^0}
\deleted{\unsure{\mathbb V}}\check{\mathbb V}(L+L^{\otimes2}).
\]
avec $\check B\otimes L$. \ldots\ Donc

\uline{Proposition} \deleted{$\unsure{\Sigma^2(F)}$} Si $F$ est un fibré affine
de rang $1$, \deleted{$\Sigma^2F$} correspondant à l'extension
\unsure{(3\ldots{}|(30\ldots{}} (37), alors $\Sigma^2(F)$ est canoniquement
muni d'une structure de fibré affine
sous le groupe vectoriel $\check B\otimes L$ (extension de
$L^{\deleted{\unsure{0}}}$ par $L^{\otimes2}$). Le fibré \deleted{associé} affine de

% archive sheet number: 30

